\section{Finite-Sample Diagnostics}\label{app:finite_sample_geometry}

This appendix gives finite-sample diagnostics for the theorem line and the
remaining operational frontier. These checks are not substitutes for theorem
statements.

\subsection{Finite-Sample Geometry Check}

For the finite-sample term, we estimate empirical $W_2$ between two i.i.d.
samples from the same distribution using exact assignment in low and moderate
dimension, then fit a log--log slope in $n$. We also probe the fixed-span
one-dimensional triangular window directly by comparing a triangular-array
sample, with one draw from each drifted component, against an i.i.d. sample from
the same window mixture. In that check the within-window span $\zeta n^H$ is
held fixed. Finally, we check the Sinkhorn proxy at a few fixed values of
$\varepsilon$.

\begin{table}[!htbp]
\centering
\scriptsize
\setlength{\tabcolsep}{4pt}
\caption{Empirical scaling check for the finite-sample term. The i.i.d. low-dimensional $W_2$ rows stay close to the expected root-$n$ regime, the one-dimensional triangular-array rows compare against the window-mixture benchmark at fixed span, and the Sinkhorn rows illustrate how the proxy changes with $\varepsilon$.}
\label{tab:finite_sample_geometry}
\resizebox{\columnwidth}{!}{%
\begin{tabular}{lrr}
\toprule
Setting & Estimated slope & Comment \\
\midrule
$W_2$, $d=1$ & -0.48 & close to $-1/2$ \\
$W_2$, $d=2$ & -0.41 & slower decay \\
$W_2$, $d=4$ & -0.25 & slower decay \\
$W_2$, $d=5$ & -0.24 & slower decay \\
$W_2$, i.i.d. mixture, $d=1$, $H=0.5$ & -0.50 & fixed span $\zeta n^H=0.25$ \\
$W_2$, triangular, $d=1$, $H=0.5$ & -0.45 & fixed span $\zeta n^H=0.25$ \\
$W_2$, i.i.d. mixture, $d=1$, $H=1.0$ & -0.47 & fixed span $\zeta n^H=0.25$ \\
$W_2$, triangular, $d=1$, $H=1.0$ & -0.41 & fixed span $\zeta n^H=0.25$ \\
Sinkhorn $\varepsilon=0.50$ & -0.21 & proxy constant changes \\
Sinkhorn $\varepsilon=0.10$ & -0.53 & proxy constant changes \\
Sinkhorn $\varepsilon=0.05$ & -0.67 & proxy constant changes \\
\bottomrule
\end{tabular}
}
\end{table}


In low dimension, the raw i.i.d. Wasserstein term is compatible with a root-$n$
rate, whereas higher-dimensional raw Wasserstein slows down. In the
one-dimensional fixed-span triangular-array check, the recovered slopes stay near
the corresponding i.i.d. mixture baseline, which is consistent with the
proof-model and design-effect reading used in the 1D proof-model proposition.

\subsection{Bahadur Remainder Diagnostics}

Two deformation families make the proof-model remainder assumption concrete. The
first is a smooth fixed-support monotone warp class
$g_\delta(x)=x+\delta x(1-x)$ with $|\delta|\le \Delta<1$. On the tested spans
$0.2$ and $0.4$, the integrated remainder decays at rates between $1.40$ and
$1.54$, and the scaled quantity $n\int_0^1 R_n(u)^2du$ still decays at rates
between $0.40$ and $0.54$. The second is the fixed-span triangular-array
uniform-kernel model used in the proof line. There the integrated residual rate
is about $1.43$ for the triangular sample and $1.49$ for the i.i.d. mixture
benchmark, while the residual-to-MSE ratio drops from about $0.38$ at $n=25$ to
about $0.11$ at $n=400$.

These diagnostics locate the role of drift span explicitly. Within the tested
fixed-span range, the remainder stays lower-order and the empirical rate remains
well above the $n^{-1}$ threshold needed for item~(iv). Rougher kernels and
growing-span experiments remain less regular, which is why the manuscript keeps
the general triangular-array verification problem open.

\subsection{Extended and Operational Regimes}

The extended rows compare triangular and i.i.d. mixture rates on embedded
low-intrinsic-dimensional support across span. The operational rows compare
fixed-$\varepsilon$ Sinkhorn rates across regularization values on embedded
low-intrinsic-dimensional support in ambient dimension eight.

\begin{table}[!htbp]
\centering
\scriptsize
\setlength{\tabcolsep}{3pt}
\caption{Representative rows for the extended and operational finite-sample regimes. The extended rows compare triangular and i.i.d. mixture rates across span, while the operational rows compare fixed-$\varepsilon$ Sinkhorn across $\varepsilon$ on embedded low-intrinsic-dimensional support.}
\label{tab:finite_sample_regimes}
\resizebox{\columnwidth}{!}{%
\begin{tabular}{llrr}
\toprule
Regime & Setting & Parameter & Rate exponent $a$ \\
\midrule
extended & triangular ambient d=8, intrinsic k=1, span=0.10 & 0.10 & 0.4652 \\
extended & iid mixture ambient d=8, intrinsic k=1, span=0.10 & 0.10 & 0.5722 \\
extended & triangular ambient d=8, intrinsic k=1, span=0.25 & 0.25 & 0.4737 \\
extended & iid mixture ambient d=8, intrinsic k=1, span=0.25 & 0.25 & 0.5588 \\
extended & triangular ambient d=8, intrinsic k=1, span=0.50 & 0.50 & 0.4792 \\
extended & iid mixture ambient d=8, intrinsic k=1, span=0.50 & 0.50 & 0.5425 \\
operational & Sinkhorn iid mixture ambient d=8, intrinsic k=1, eps=0.50 & 0.50 & 0.5599 \\
operational & Sinkhorn triangular ambient d=8, intrinsic k=1, eps=0.50 & 0.50 & 0.5706 \\
operational & Sinkhorn iid mixture ambient d=8, intrinsic k=1, eps=0.20 & 0.20 & 0.5466 \\
operational & Sinkhorn triangular ambient d=8, intrinsic k=1, eps=0.20 & 0.20 & 0.6150 \\
operational & Sinkhorn iid mixture ambient d=8, intrinsic k=1, eps=0.10 & 0.10 & 0.4646 \\
operational & Sinkhorn triangular ambient d=8, intrinsic k=1, eps=0.10 & 0.10 & 0.5060 \\
operational & Sinkhorn iid mixture ambient d=8, intrinsic k=1, eps=0.05 & 0.05 & 0.4339 \\
operational & Sinkhorn triangular ambient d=8, intrinsic k=1, eps=0.05 & 0.05 & 0.5124 \\
\bottomrule
\end{tabular}
}
\end{table}


The extended-regime rows indicate that low intrinsic dimension preserves a rate
near the root-$n$ regime without a dramatic triangular-versus-benchmark gap. The
operational-regime rows record empirical evidence relevant to the fixed-
$\varepsilon$ Sinkhorn setting in \cref{def:fixed_epsilon_sinkhorn_setting}:
together with exact support-complexity inheritance and the dual threshold
$\alpha>k/2$, they align with the embedded bandwise statement. A full proof of
triangular-array inheritance would close this loop.

Two complementary diagnostics sharpen the same boundary. First, direct probes of
the centered Sinkhorn scaling Jacobian on the embedded fixed-span model show
that bands approaching zero are the numerically unstable region, whereas on a
moderate band such as $[0.2,0.8]$ the cross-coupling block remains contractive
and the self-coupling inverse norms appear to plateau with sample size. Second,
on the supported synthetic frontier grid, intrinsic-dimension estimates
predict the observed stable-band cutoff better than ambient dimension.
These diagnostics identify a useful regime proxy and isolate the band where the
remaining instability is concentrated. They do not upgrade the embedded fixed-
$\varepsilon$ statement from conjecture to theorem.
